% Options for packages loaded elsewhere
\PassOptionsToPackage{unicode}{hyperref}
\PassOptionsToPackage{hyphens}{url}
%
\documentclass[
  11pt,
]{article}
\usepackage{amsmath,amssymb}
\usepackage{iftex}
\ifPDFTeX
  \usepackage[T1]{fontenc}
  \usepackage[utf8]{inputenc}
  \usepackage{textcomp} % provide euro and other symbols
\else % if luatex or xetex
  \usepackage{unicode-math} % this also loads fontspec
  \defaultfontfeatures{Scale=MatchLowercase}
  \defaultfontfeatures[\rmfamily]{Ligatures=TeX,Scale=1}
\fi
\usepackage{lmodern}
\ifPDFTeX\else
  % xetex/luatex font selection
\fi
% Use upquote if available, for straight quotes in verbatim environments
\IfFileExists{upquote.sty}{\usepackage{upquote}}{}
\IfFileExists{microtype.sty}{% use microtype if available
  \usepackage[]{microtype}
  \UseMicrotypeSet[protrusion]{basicmath} % disable protrusion for tt fonts
}{}
\makeatletter
\@ifundefined{KOMAClassName}{% if non-KOMA class
  \IfFileExists{parskip.sty}{%
    \usepackage{parskip}
  }{% else
    \setlength{\parindent}{0pt}
    \setlength{\parskip}{6pt plus 2pt minus 1pt}}
}{% if KOMA class
  \KOMAoptions{parskip=half}}
\makeatother
\usepackage{xcolor}
\usepackage[margin=2.2cm]{geometry}
\usepackage{longtable,booktabs,array}
\usepackage{calc} % for calculating minipage widths
% Correct order of tables after \paragraph or \subparagraph
\usepackage{etoolbox}
\makeatletter
\patchcmd\longtable{\par}{\if@noskipsec\mbox{}\fi\par}{}{}
\makeatother
% Allow footnotes in longtable head/foot
\IfFileExists{footnotehyper.sty}{\usepackage{footnotehyper}}{\usepackage{footnote}}
\makesavenoteenv{longtable}
\setlength{\emergencystretch}{3em} % prevent overfull lines
\providecommand{\tightlist}{%
  \setlength{\itemsep}{0pt}\setlength{\parskip}{0pt}}
\setcounter{secnumdepth}{-\maxdimen} % remove section numbering
\ifLuaTeX
  \usepackage{selnolig}  % disable illegal ligatures
\fi
\IfFileExists{bookmark.sty}{\usepackage{bookmark}}{\usepackage{hyperref}}
\IfFileExists{xurl.sty}{\usepackage{xurl}}{} % add URL line breaks if available
\urlstyle{same}
\hypersetup{
  pdftitle={Worksheet 4B: What Derivatives and Integrals Tell Us},
  hidelinks,
  pdfcreator={LaTeX via pandoc}}

\title{Worksheet 4B: What Derivatives and Integrals Tell Us}
\author{}
\date{}

\begin{document}
\maketitle

\begin{quote}
\textbf{What This Worksheet Is About}

We've learned \emph{how} to compute derivatives and integrals. Now:
\emph{what do they mean}?

\begin{itemize}
\tightlist
\item
  \textbf{Derivatives} → rate of change, slope, velocity,
  increasing/decreasing
\item
  \textbf{Integrals} → accumulation, area, total distance, total
  quantity
\end{itemize}

These interpretations are why calculus is so powerful!
\end{quote}

\begin{center}\rule{0.5\linewidth}{0.5pt}\end{center}

\hypertarget{part-a-derivatives-rate-of-change-and-slope-16-problems}{%
\subsection{Part A: Derivatives --- Rate of Change and Slope (16
problems)}\label{part-a-derivatives-rate-of-change-and-slope-16-problems}}

\begin{longtable}[]{@{}ll@{}}
\toprule\noalign{}
Derivative & Meaning \\
\midrule\noalign{}
\endhead
\bottomrule\noalign{}
\endlastfoot
\(f'(x) > 0\) & Function increasing ↗ \\
\(f'(x) < 0\) & Function decreasing ↘ \\
\(f'(x) = 0\) & Horizontal tangent (potential max/min) \\
\(f''(x) > 0\) & Concave up ∪ \\
\(f''(x) < 0\) & Concave down ∩ \\
\end{longtable}

\begin{center}\rule{0.5\linewidth}{0.5pt}\end{center}

\textbf{1.} For \(f(x) = x^2\), \(f'(x) = 2x\).

\begin{itemize}
\item
  \begin{enumerate}
  \def\labelenumi{(\alph{enumi})}
  \tightlist
  \item
    Where is \(f\) increasing? Decreasing?
  \end{enumerate}
\item
  \begin{enumerate}
  \def\labelenumi{(\alph{enumi})}
  \setcounter{enumi}{1}
  \tightlist
  \item
    Where is the tangent horizontal?
  \end{enumerate}
\item
  \begin{enumerate}
  \def\labelenumi{(\alph{enumi})}
  \setcounter{enumi}{2}
  \tightlist
  \item
    What is the slope at \(x = 3\)?
  \end{enumerate}
\end{itemize}

\textbf{2.} For \(g(x) = x^3 - 3x\):

\begin{itemize}
\item
  \begin{enumerate}
  \def\labelenumi{(\alph{enumi})}
  \tightlist
  \item
    Find \(g'(x)\).
  \end{enumerate}
\item
  \begin{enumerate}
  \def\labelenumi{(\alph{enumi})}
  \setcounter{enumi}{1}
  \tightlist
  \item
    Find critical points (where \(g'(x) = 0\)).
  \end{enumerate}
\item
  \begin{enumerate}
  \def\labelenumi{(\alph{enumi})}
  \setcounter{enumi}{2}
  \tightlist
  \item
    Classify as local max or min.
  \end{enumerate}
\end{itemize}

\textbf{3.} For \(h(t) = t^4 - 8t^2\):

\begin{itemize}
\item
  \begin{enumerate}
  \def\labelenumi{(\alph{enumi})}
  \tightlist
  \item
    Find all critical points.
  \end{enumerate}
\item
  \begin{enumerate}
  \def\labelenumi{(\alph{enumi})}
  \setcounter{enumi}{1}
  \tightlist
  \item
    Use the first or second derivative test to classify each.
  \end{enumerate}
\end{itemize}

\textbf{4.} Find the tangent line to \(f(x) = x^2\) at \((2, 4)\).

\textbf{5.} Find the tangent line to \(g(x) = x^3 - x\) at \(x = 1\).

\textbf{6.} At what point(s) on \(y = x^3 - 12x\) is the tangent
horizontal?

\begin{center}\rule{0.5\linewidth}{0.5pt}\end{center}

\textbf{Physics Applications:}

\textbf{7.} Height: \(h(t) = 20t - 5t^2\) meters.

\begin{itemize}
\item
  \begin{enumerate}
  \def\labelenumi{(\alph{enumi})}
  \tightlist
  \item
    Find velocity \(v(t) = h'(t)\).
  \end{enumerate}
\item
  \begin{enumerate}
  \def\labelenumi{(\alph{enumi})}
  \setcounter{enumi}{1}
  \tightlist
  \item
    When is velocity zero? What's happening then?
  \end{enumerate}
\item
  \begin{enumerate}
  \def\labelenumi{(\alph{enumi})}
  \setcounter{enumi}{2}
  \tightlist
  \item
    Find acceleration \(a(t) = h''(t)\). Interpret it.
  \end{enumerate}
\end{itemize}

\textbf{8.} Position: \(s(t) = t^3 - 6t^2 + 9t\) meters.

\begin{itemize}
\item
  \begin{enumerate}
  \def\labelenumi{(\alph{enumi})}
  \tightlist
  \item
    Find velocity and acceleration.
  \end{enumerate}
\item
  \begin{enumerate}
  \def\labelenumi{(\alph{enumi})}
  \setcounter{enumi}{1}
  \tightlist
  \item
    When is the object at rest?
  \end{enumerate}
\item
  \begin{enumerate}
  \def\labelenumi{(\alph{enumi})}
  \setcounter{enumi}{2}
  \tightlist
  \item
    When is acceleration zero?
  \end{enumerate}
\end{itemize}

\begin{center}\rule{0.5\linewidth}{0.5pt}\end{center}

\textbf{Second Derivative and Concavity:}

\textbf{9.} For \(f(x) = x^3\): - Where is \(f''(x) > 0\)? (Concave up)
- Where is \(f''(x) < 0\)? (Concave down) - Where is the inflection
point?

\textbf{10.} For \(g(x) = x^4 - 6x^2\): - Find inflection points. -
Describe concavity on each interval.

\begin{center}\rule{0.5\linewidth}{0.5pt}\end{center}

\textbf{Optimization:}

\textbf{11.} Profit: \(P(x) = -x^2 + 80x - 1200\). - Find the quantity
maximizing profit. - Find maximum profit.

\textbf{12.} Area with fixed perimeter: A rectangle has perimeter 100m.
- Express area as \(A(w) = w(50 - w) = 50w - w^2\). - Find dimensions
maximizing area.

\textbf{13.} A ball's height is \(h(t) = -5t^2 + 30t + 10\). - When does
it reach maximum height? - What is the maximum height?

\textbf{14.} Revenue: \(R(p) = 500p - 2p^2\). - Find price maximizing
revenue. - Find maximum revenue.

\textbf{15.} Cost: \(C(q) = q^3 - 12q^2 + 60q + 100\). - Find marginal
cost \(C'(q)\). - At what quantity is marginal cost minimized?

\textbf{16.} Summary: The derivative being zero tells us
\_\_\_\_\_\_\_\_\_\_, and the second derivative's sign tells us
\_\_\_\_\_\_\_\_\_\_.

\begin{center}\rule{0.5\linewidth}{0.5pt}\end{center}

\hypertarget{part-b-integrals-accumulation-and-area-16-problems}{%
\subsection{Part B: Integrals --- Accumulation and Area (16
problems)}\label{part-b-integrals-accumulation-and-area-16-problems}}

\begin{longtable}[]{@{}ll@{}}
\toprule\noalign{}
Integral & Meaning \\
\midrule\noalign{}
\endhead
\bottomrule\noalign{}
\endlastfoot
\(\int_a^b f(x) \, dx\) & Signed area under curve \\
\(\int v(t) \, dt\) & Displacement \\
\(\int r(t) \, dt\) & Total quantity accumulated \\
\end{longtable}

\begin{center}\rule{0.5\linewidth}{0.5pt}\end{center}

\textbf{Area Under Curves:}

\textbf{17.} Find the area under \(f(x) = 3\) from \(x = 1\) to
\(x = 5\). - Geometrically: rectangle with base \_\_\_ and height \_\_\_
- By integral: \(\int_1^5 3 \, dx ={}\)

\textbf{18.} Find the area under \(f(x) = x\) from \(x = 0\) to
\(x = 4\). - Geometrically: triangle - By integral:
\(\int_0^4 x \, dx ={}\)

\textbf{19.} \(\int_0^2 x^2 \, dx ={}\)

\textbf{20.} \(\int_1^3 (2x + 1) \, dx ={}\)

\begin{center}\rule{0.5\linewidth}{0.5pt}\end{center}

\textbf{Signed Area:}

\textbf{21.} \(\int_{-2}^2 x \, dx ={}\) \_\_\_\_\_

Why? (Think about symmetry)

\textbf{22.} \(\int_{-1}^1 x^3 \, dx ={}\) \_\_\_\_\_

This is an odd function on a symmetric interval.

\textbf{23.} \(\int_{-2}^2 x^2 \, dx ={}\) \_\_\_\_\_

This is an even function: \(\int_{-a}^a f(x) dx = 2\int_0^a f(x) dx\)

\textbf{24.} For \(f(x) = x - 1\) from \(x = 0\) to \(x = 2\): - Sketch
the function. - Where is it positive? Negative? - Calculate
\(\int_0^2 (x - 1) dx\). Interpret the result.

\begin{center}\rule{0.5\linewidth}{0.5pt}\end{center}

\textbf{Physical Applications:}

\textbf{25.} Velocity: \(v(t) = 3t^2\) m/s. Find displacement from
\(t = 0\) to \(t = 2\).

\textbf{26.} Velocity: \(v(t) = t - 2\) m/s from \(t = 0\) to \(t = 4\).
- (a) When is the object moving forward? Backward? - (b) Find net
displacement: \(\int_0^4 (t - 2) dt\) - (c) Find total distance
traveled.

\textbf{27.} Flow rate: \(r(t) = 100 - 4t\) liters/min. - Find total
flow from \(t = 0\) to \(t = 10\).

\textbf{28.} Population growth rate: \(P'(t) = 200 + 10t\) people/year.
- Find total growth over 5 years.

\begin{center}\rule{0.5\linewidth}{0.5pt}\end{center}

\textbf{Accumulation Functions:}

\textbf{29.} Let \(A(x) = \int_0^x t^2 \, dt\). - (a) Compute \(A(x)\)
explicitly. - (b) Find \(A(1)\), \(A(2)\), \(A(3)\). - (c) Find
\(A'(x)\). What do you notice?

\textbf{30.} If \(F(x) = \int_1^x (t^3 + 1) dt\), what is \(F'(x)\)?

(Use the Fundamental Theorem --- no need to compute the integral!)

\textbf{31.} If \(G(x) = \int_0^x \sqrt{t} \, dt\), what is \(G'(x)\)?

\textbf{32.} Explain: ``Integrating then differentiating gets you back
to where you started.''

\begin{center}\rule{0.5\linewidth}{0.5pt}\end{center}

\hypertarget{part-c-the-fundamental-theorem-tying-it-together-12-problems}{%
\subsection{Part C: The Fundamental Theorem --- Tying It Together (12
problems)}\label{part-c-the-fundamental-theorem-tying-it-together-12-problems}}

\textbf{The Fundamental Theorem of Calculus:}

\begin{enumerate}
\def\labelenumi{\arabic{enumi}.}
\item
  \(\frac{d}{dx}\left[\int_a^x f(t) \, dt\right] = f(x)\)
\item
  \(\int_a^b f(x) \, dx = F(b) - F(a)\) where \(F' = f\)
\end{enumerate}

\begin{center}\rule{0.5\linewidth}{0.5pt}\end{center}

\textbf{Using Part 1:}

\textbf{33.} \(\frac{d}{dx}\left[\int_0^x t^4 \, dt\right] ={}\)

\textbf{34.} \(\frac{d}{dx}\left[\int_2^x (3s^2 - s) \, ds\right] ={}\)

\textbf{35.}
\(\frac{d}{dx}\left[\int_1^x \frac{1}{t^2} \, dt\right] ={}\)

\begin{center}\rule{0.5\linewidth}{0.5pt}\end{center}

\textbf{Using Part 2:}

\textbf{36.} \(\int_0^3 2x \, dx ={}\) (use antiderivative \(x^2\))

\textbf{37.} \(\int_1^4 (x^2 - 1) \, dx ={}\)

\textbf{38.} \(\int_0^8 x^{1/3} \, dx ={}\)

\textbf{39.} \(\int_1^e \frac{1}{x} \, dx ={}\) (antiderivative is
\(\ln x\); we'll derive this in 4D!)

\begin{center}\rule{0.5\linewidth}{0.5pt}\end{center}

\textbf{Round-Trip Verification:}

\textbf{40.} Start with \(f(x) = x^3\). - Integrate:
\(\int f(x) dx ={}\) - Differentiate your answer. - Did you get \(f(x)\)
back?

\textbf{41.} Start with \(g(x) = 6x^2 - 4x\). - Differentiate:
\(g'(x) ={}\) - Integrate your answer. - Did you get \(g(x)\) back (plus
C)?

\textbf{42.} Explain the asymmetry: - Differentiating then integrating
gives \(f(x) + C\) - Integrating then differentiating gives exactly
\(f(x)\)

Why the difference?

\begin{center}\rule{0.5\linewidth}{0.5pt}\end{center}

\textbf{Connecting Position, Velocity, Acceleration:}

\textbf{43.} Complete the chain:

Position \(s(t)\) \(\xrightarrow{\text{derivative}}\) \_\_\_\_\_\_\_\_
\(\xrightarrow{\text{derivative}}\) \_\_\_\_\_\_\_\_

\_\_\_\_\_\_\_\_ \(\xrightarrow{\text{integral}}\) Velocity
\(\xrightarrow{\text{integral}}\) \_\_\_\_\_\_\_\_

\textbf{44.} An object has acceleration \(a(t) = 6t\), with \(v(0) = 2\)
and \(s(0) = 1\). - Find \(v(t)\) by integrating \(a(t)\) and using the
initial condition. - Find \(s(t)\) by integrating \(v(t)\) and using the
initial condition.

\begin{center}\rule{0.5\linewidth}{0.5pt}\end{center}

\hypertarget{part-d-area-between-curves-8-problems}{%
\subsection{Part D: Area Between Curves (8
problems)}\label{part-d-area-between-curves-8-problems}}

Area between \(f(x)\) and \(g(x)\) from \(a\) to \(b\):

\[\text{Area} = \int_a^b |f(x) - g(x)| \, dx\]

If \(f(x) \geq g(x)\) on \([a,b]\): Area
\(= \int_a^b (f(x) - g(x)) \, dx\)

\begin{center}\rule{0.5\linewidth}{0.5pt}\end{center}

\textbf{45.} Find area between \(y = x^2\) and \(y = x\) from \(x = 0\)
to \(x = 1\). - Which is on top? - Set up and compute the integral.

\textbf{46.} Find area between \(y = x^2\) and \(y = 4\). - Find
intersection points. - Set up and compute.

\textbf{47.} Find area enclosed by \(y = x^2\) and \(y = 2x\). - Find
intersections. - Integrate from intersection to intersection.

\textbf{48.} Find area between \(y = x^3\) and \(y = x\) from \(x = -1\)
to \(x = 1\). - Careful: which is on top changes at \(x = 0\)!

\begin{center}\rule{0.5\linewidth}{0.5pt}\end{center}

\textbf{Applications:}

\textbf{49.} Consumer Surplus: Demand is \(p = 50 - 2q\), equilibrium
price is \(p^* = 30\). - Consumer surplus =
\(\int_0^{q^*} (50 - 2q - 30) dq\) - Find \(q^*\) and calculate surplus.

\textbf{50.} The area between a velocity curve and the t-axis represents
\_\_\_\_\_\_\_\_\_\_.

\textbf{51.} The area between two position curves represents the
\_\_\_\_\_\_\_\_\_\_ between them.

\textbf{52.} Gini Coefficient: Perfect equality is \(y = x\). A Lorenz
curve is \(y = x^2\). - Gini =
\(\frac{\text{Area between curves}}{\text{Area under } y = x}\) -
Calculate it.

\begin{center}\rule{0.5\linewidth}{0.5pt}\end{center}

\hypertarget{part-e-connecting-to-ml-and-optimization-8-problems}{%
\subsection{Part E: Connecting to ML and Optimization (8
problems)}\label{part-e-connecting-to-ml-and-optimization-8-problems}}

\textbf{53.} Loss function: \(L(w) = w^2 - 6w + 10\). - Find \(L'(w)\).
- Find the minimum. - Verify with second derivative.

\textbf{54.} Gradient descent update:
\(w_{new} = w_{old} - \alpha \cdot L'(w_{old})\)

For \(L(w) = w^2 - 4w + 5\) with \(\alpha = 0.1\), starting at
\(w_0 = 0\): - Calculate \(w_1\), \(w_2\), \(w_3\). - What value is
\(w\) approaching?

\textbf{55.} Total loss over dataset \([0, 4]\) with loss function
\(\ell(x) = (x - 2)^2\):

Total = \(\int_0^4 (x-2)^2 dx\). Calculate it.

\textbf{56.} Mean loss = \(\frac{1}{4}\int_0^4 (x-2)^2 dx\). Calculate
it.

\textbf{57.} Regularized loss: \(L(w) = (w - 5)^2 + 0.5w^2\). - Expand
and simplify. - Find optimal \(w\). - Compare to unregularized optimum
\(w = 5\).

\textbf{58.} Explain: Why is finding where \(L'(w) = 0\) central to
training ML models?

\textbf{59.} Explain: Why might we care about \(\int L(w) dw\) over a
region of weight space?

\textbf{60.} The Fundamental Theorem connects local information
(derivative at a point) to global information (integral over an
interval). How does this relate to: - Local vs.~global minima in
optimization? - Gradient (local) vs.~total loss (global)?

\begin{center}\rule{0.5\linewidth}{0.5pt}\end{center}

\end{document}
